% Supplementary material of the xg-vendor paper draft (xg-vendor 17). Every number is \val{id} from numbers.tex.
\documentclass[10pt,twocolumn,letterpaper]{article}

\usepackage[review]{cvpr}

\definecolor{cvprblue}{rgb}{0.21,0.49,0.74}
\usepackage[pagebackref,breaklinks,colorlinks,allcolors=cvprblue]{hyperref}
\usepackage{tabularx}

% Generated by scripts/check_vendor_paper.py --write from claims.toml; do not edit.
% \val{id} prints the number of claim id; an unknown id stops the LaTeX run.
\providecommand{\defval}[2]{\expandafter\def\csname xgv:#1\endcsname{#2}}
\providecommand{\val}{}
\DeclareRobustCommand{\val}[1]{\ifcsname xgv:#1\endcsname\csname xgv:#1\endcsname\else\PackageError{numbers}{No number with id #1}{}\fi}
\defval{pg.code_files}{32}
\defval{run.gate_thr}{0.005}
\defval{run.prereg_sha}{312e4e4a3282}
\defval{pg.nf_matches}{86}
\defval{pg.recovered_pen}{7}
\defval{pg.ci95}{95\%}
\defval{pg.corner_margin}{0.021}
\defval{pg.near}{0.02}
\defval{run.delta}{0.005}
\defval{run.n_boot}{5,000}
\defval{pg.alpha}{5\%}
\defval{pg.ci90}{90\%}
\defval{pg.bonf}{2.5\%}
\defval{pa.p_train}{1,478}
\defval{pa.p_valid}{639}
\defval{pa.p_test}{637}
\defval{pa.secondary}{797}
\defval{pa.recovered}{50}
\defval{rd.cut}{209}
\defval{pa.sb_unpaired}{875}
\defval{pa.mispairs}{26}
\defval{pa.mispair_s}{10 s}
\defval{ti.unpaired}{285}
\defval{ti.r_shot}{36}
\defval{pa.no_sb_10s}{249}
\defval{pa.window}{2.5 s}
\defval{ti.r_other}{176}
\defval{ti.r_dead}{52}
\defval{ti.r_none}{21}
\defval{ti.verdict}{``independent''}
\defval{ti.n_subs}{686}
\defval{ti.min_subs}{50}
\defval{ti.agree_ms}{\ensuremath{\pm}200 ms}
\defval{ti.subs_share}{13.3\%}
\defval{ti.shared}{90\%}
\defval{ti.subs_median}{1788 ms}
\defval{ti.lag}{1000 ms}
\defval{ti.shots}{96.6\%}
\defval{ti.kick}{80\%}
\defval{ti.corners}{82.1\%}
\defval{pa.paired}{2,754}
\defval{pa.matches}{135}
\defval{pa.index_shots}{3,039}
\defval{pa.sn_matches}{500}
\defval{pa.rate}{90.62\%}
\defval{pa.k0_rate}{85\%}
\defval{pa.goal_rate}{99.53\%}
\defval{pa.k0_goal}{99\%}
\defval{pa.blocked}{515/890}
\defval{pa.wayward}{47/146}
\defval{pa.dropped}{74}
\defval{co.pool_shots}{36,626}
\defval{co.pool_matches}{1,491}
\defval{co.k1_ff}{0.24607}
\defval{co.k1_vendor}{0.24732}
\defval{co.k1_gap}{\ensuremath{-}0.0012}
\defval{co.k1_margin}{0.01}
\defval{co.k1_pass}{passed}
\defval{sp.eval_train}{72}
\defval{sp.eval_valid}{31}
\defval{gc.matches}{37}
\defval{gc.pass}{passed}
\defval{gc.ll_new}{0.32655}
\defval{gc.ll_old}{0.32380}
\defval{gc.ll_diff}{\ensuremath{+}0.00275 [\ensuremath{-}0.00096, \ensuremath{+}0.00650]}
\defval{gc.tol_ll}{\ensuremath{+}0.010}
\defval{gc.auc_new}{0.79240}
\defval{gc.auc_old}{0.79884}
\defval{gc.auc_diff}{\ensuremath{-}0.00643 [\ensuremath{-}0.01721, \ensuremath{+}0.00465]}
\defval{gc.tol_auc}{\ensuremath{-}0.020}
\defval{gc.shots}{825}
\defval{sp.training}{328}
\defval{cf.s1.p.chosen}{2,754}
\defval{cf.s1.p.label}{0}
\defval{cf.s1.p.pen}{33}
\defval{cf.s1.p.fx}{0}
\defval{cf.s1.p.nf}{0}
\defval{cf.s1.p.n}{2,721}
\defval{cf.s1.p.goals}{374}
\defval{cf.s1.p.matches}{135}
\defval{cf.s1.c.chosen}{3,601}
\defval{cf.s1.c.label}{12}
\defval{cf.s1.c.pen}{40}
\defval{cf.s1.c.fx}{1}
\defval{cf.s1.c.nf}{147}
\defval{cf.s1.c.n}{3,401}
\defval{cf.s1.c.goals}{417}
\defval{cf.s1.c.matches}{135}
\defval{cf.c.p.frame_vendor}{\ensuremath{+}0.0283 [\ensuremath{+}0.0168, \ensuremath{+}0.0392]}
\defval{cf.c.c.frame_vendor}{\ensuremath{+}0.0282 [\ensuremath{+}0.0184, \ensuremath{+}0.0377]}
\defval{cf.c.p.frame_ff}{\ensuremath{+}0.0272 [\ensuremath{+}0.0162, \ensuremath{+}0.0380]}
\defval{cf.c.c.frame_ff}{\ensuremath{+}0.0264 [\ensuremath{+}0.0167, \ensuremath{+}0.0355]}
\defval{cf.c.p.frame_da}{\ensuremath{+}0.0029 [\ensuremath{-}0.0083, \ensuremath{+}0.0137]}
\defval{cf.c.c.frame_da}{\ensuremath{-}0.0012 [\ensuremath{-}0.0105, \ensuremath{+}0.0075]}
\defval{cf.c.p.vendor_ff}{\ensuremath{-}0.0010 [\ensuremath{-}0.0076, \ensuremath{+}0.0057]}
\defval{cf.c.c.vendor_ff}{\ensuremath{-}0.0018 [\ensuremath{-}0.0073, \ensuremath{+}0.0037]}
\defval{cf.c.p.ff_event}{\ensuremath{-}0.0059 [\ensuremath{-}0.0109, \ensuremath{-}0.0009]}
\defval{cf.c.c.ff_event}{\ensuremath{-}0.0058 [\ensuremath{-}0.0101, \ensuremath{-}0.0013]}
\defval{cf.c.p.event_da}{\ensuremath{-}0.0184 [\ensuremath{-}0.0272, \ensuremath{-}0.0101]}
\defval{cf.c.c.event_da}{\ensuremath{-}0.0219 [\ensuremath{-}0.0300, \ensuremath{-}0.0146]}
\defval{cf.c.p.da_const}{\ensuremath{-}0.0488 [\ensuremath{-}0.0608, \ensuremath{-}0.0376]}
\defval{cf.c.c.da_const}{\ensuremath{-}0.0413 [\ensuremath{-}0.0506, \ensuremath{-}0.0323]}
\defval{cf.h3.f0.angle}{0.450 [0.409, 0.489]}
\defval{cf.h3.f0.ff_def_cone}{0.367 [0.319, 0.409]}
\defval{cf.h3.f0.ge2}{0.849 [0.826, 0.871]}
\defval{cf.h3.f0.gk_depth}{0.325 [0.269, 0.377]}
\defval{cf.h3.f0.gk_off}{0.025 [0.004, 0.063]}
\defval{cf.h3.f0.gk_dist}{0.318 [0.244, 0.384]}
\defval{cf.h3.f0.opp}{0.185 [0.149, 0.218]}
\defval{cf.h3.f0.head}{0.906 [0.883, 0.927]}
\defval{cf.h3.thr.angle}{0.25}
\defval{cf.h3.thr.head}{0.70}
\defval{cf.h4.slope}{\ensuremath{+}0.545 [\ensuremath{+}0.259, \ensuremath{+}0.857]}
\defval{cf.h4.slope_200}{\ensuremath{-}0.212 [\ensuremath{-}0.822, \ensuremath{+}0.388]}
\defval{cf.h4.early}{\ensuremath{-}0.073 [\ensuremath{-}0.150, \ensuremath{+}0.007]}
\defval{cf.h4.early_200}{\ensuremath{+}0.005 [\ensuremath{-}0.072, \ensuremath{+}0.081]}
\defval{cf.h4.late}{\ensuremath{-}0.047 [\ensuremath{-}0.114, \ensuremath{+}0.020]}
\defval{cf.h4.late_200}{\ensuremath{-}0.058 [\ensuremath{-}0.128, \ensuremath{+}0.012]}
\defval{cf.h4.n}{2,347}
\defval{cf.h4.n200}{2,284}
\defval{cf.h4.off_on}{\ensuremath{+}0.025 [\ensuremath{-}0.009, \ensuremath{+}0.057]}
\defval{cf.h4.off_on_u}{\ensuremath{+}0.041 [\ensuremath{+}0.007, \ensuremath{+}0.073]}
\defval{cf.pen.p.n}{33}
\defval{cf.pen.p.goals}{22}
\defval{cf.pen.p.vp}{0.783}
\defval{cf.pen.p.vll}{0.6727}
\defval{cf.pen.p.fp}{0.697}
\defval{cf.pen.p.fll}{0.6387}
\defval{cf.pen.p.pp}{0.758}
\defval{cf.pen.p.pll}{0.6578}
\defval{cf.pen.c.n}{40}
\defval{cf.pen.c.goals}{29}
\defval{cf.pen.c.vp}{0.783}
\defval{cf.pen.c.vll}{0.5977}
\defval{cf.pen.c.fp}{0.697}
\defval{cf.pen.c.fll}{0.5900}
\defval{cf.pen.c.pp}{0.758}
\defval{cf.pen.c.pll}{0.5911}
\defval{cf.h2.p.diff}{\ensuremath{-}0.0004 [\ensuremath{-}0.0026, \ensuremath{+}0.0018]}
\defval{cf.h2.p.tost}{0.0006}
\defval{cf.h3.f0.dist}{0.711 [0.678, 0.739]}
\defval{cf.h3.f0.cone}{0.761 [0.743, 0.778]}
\defval{cf.m.p.constant}{0.4006 [0.3751, 0.4271]}
\defval{cf.a.p.constant}{0.477}
\defval{cf.m.c.constant}{0.3722}
\defval{cf.m.p.dist_angle}{0.3518 [0.3257, 0.3783]}
\defval{cf.a.p.dist_angle}{0.741}
\defval{cf.m.c.dist_angle}{0.3309}
\defval{cf.m.p.event}{0.3333 [0.3076, 0.3591]}
\defval{cf.a.p.event}{0.778}
\defval{cf.m.c.event}{0.3091}
\defval{cf.m.p.freeze_frame}{0.3275 [0.3016, 0.3538]}
\defval{cf.a.p.freeze_frame}{0.788}
\defval{cf.m.c.freeze_frame}{0.3033}
\defval{cf.m.p.vendor}{0.3264 [0.3016, 0.3515]}
\defval{cf.a.p.vendor}{0.792}
\defval{cf.m.c.vendor}{0.3015}
\defval{cf.m.p.frame}{0.3547 [0.3298, 0.3802]}
\defval{cf.a.p.frame}{0.736}
\defval{cf.m.c.frame}{0.3297}
\defval{cf.m.p.frame_ff}{0.3271 [0.3017, 0.3538]}
\defval{cf.a.p.frame_ff}{0.789}
\defval{cf.m.c.frame_ff}{0.3028}
\defval{cf.m.p.frame_vendor}{0.3258 [0.3016, 0.3510]}
\defval{cf.a.p.frame_vendor}{0.792}
\defval{cf.m.c.frame_vendor}{0.3007}
\defval{cf.c.p.auc_frame_vendor}{\ensuremath{-}0.056 [\ensuremath{-}0.080, \ensuremath{-}0.031]}
\defval{cf.r.p.vendor}{\ensuremath{-}0.0328 [\ensuremath{-}0.0456, \ensuremath{-}0.0205]}
\defval{cf.r.p.ff}{\ensuremath{-}0.0274 [\ensuremath{-}0.0399, \ensuremath{-}0.0155]}
\defval{cf.r.p.frame}{\ensuremath{-}0.0041 [\ensuremath{-}0.0182, \ensuremath{+}0.0092]}
\defval{cf.h2.p.ci90}{[\ensuremath{-}0.0022, \ensuremath{+}0.0015]}
\defval{cf.h2.p.sup}{0.5274}
\defval{cf.h2.p.margin}{0.0022}
\defval{cf.h2.c.diff}{\ensuremath{-}0.0005 [\ensuremath{-}0.0025, \ensuremath{+}0.0012]}
\defval{cf.h2.c.ci90}{[\ensuremath{-}0.0021, \ensuremath{+}0.0010]}
\defval{cf.h2.c.tost}{0.0000}
\defval{cf.h2.c.sup}{0.4030}
\defval{cf.h2.c.margin}{0.0021}
\defval{cf.fv.p.diff}{\ensuremath{-}0.0007 [\ensuremath{-}0.0034, \ensuremath{+}0.0020]}
\defval{cf.fv.p.ci90}{[\ensuremath{-}0.0028, \ensuremath{+}0.0016]}
\defval{cf.fv.p.tost}{0.0042}
\defval{cf.fv.p.sup}{0.4384}
\defval{cf.fv.p.margin}{0.0029}
\defval{cf.fv.c.diff}{\ensuremath{-}0.0008 [\ensuremath{-}0.0032, \ensuremath{+}0.0012]}
\defval{cf.fv.c.ci90}{[\ensuremath{-}0.0027, \ensuremath{+}0.0010]}
\defval{cf.fv.c.tost}{0.0018}
\defval{cf.fv.c.sup}{0.3158}
\defval{cf.fv.c.margin}{0.0027}
\defval{cf.h2.p.reading}{``equivalent at $\delta$: the frame adds no more than $\delta$''}
\defval{cf.h2.p.point}{\ensuremath{-}0.0004}
\defval{cf.h2.p.ci95}{[\ensuremath{-}0.0026, \ensuremath{+}0.0018]}
\defval{cf.h2.c.reading}{equivalent at $\delta$: the frame adds no more than $\delta$}
\defval{cf.h2.p.sd}{0.00111}
\defval{cf.ctl.f0}{\ensuremath{-}0.0006}
\defval{cf.ctl.f1}{\ensuremath{+}0.0005}
\defval{cf.ctl.dino}{\ensuremath{-}0.0002}
\defval{cf.h3.n}{2,721}
\defval{cf.h3.thr.dist}{0.50}
\defval{cf.h3.f0.dist.reading}{encodes}
\defval{cf.h3.f1.dist}{0.716 [0.685, 0.743]}
\defval{cf.h3.f1.dist.reading}{encodes}
\defval{cf.h3.d.dist}{0.684 [0.656, 0.711]}
\defval{cf.d.f0.dist}{\ensuremath{+}0.027 [\ensuremath{+}0.011, \ensuremath{+}0.042]}
\defval{cf.h3.thr.cone}{0.70}
\defval{cf.h3.f0.cone.reading}{encodes}
\defval{cf.h3.f1.cone}{0.758 [0.741, 0.775]}
\defval{cf.h3.f1.cone.reading}{encodes}
\defval{cf.h3.d.cone}{0.735 [0.716, 0.754]}
\defval{cf.d.f0.cone}{\ensuremath{+}0.025 [\ensuremath{+}0.013, \ensuremath{+}0.038]}
\defval{cf.d.f1.dist}{\ensuremath{+}0.032 [\ensuremath{+}0.018, \ensuremath{+}0.046]}
\defval{cf.d.f1.cone}{\ensuremath{+}0.023 [\ensuremath{+}0.010, \ensuremath{+}0.035]}
\defval{cf.h4.on}{1,079}
\defval{cf.h4.main}{\ensuremath{+}0.056 [\ensuremath{+}0.004, \ensuremath{+}0.108]}
\defval{cf.h4.main_u}{\ensuremath{-}0.034 [\ensuremath{-}0.107, \ensuremath{+}0.036]}
\defval{cf.h4.off_mean}{\ensuremath{+}0.057 [\ensuremath{+}0.041, \ensuremath{+}0.073]}
\defval{dv.target}{0.80}
\defval{dv.pow002}{0.284}
\defval{dv.pow003}{0.509}
\defval{dv.pow005}{0.849}
\defval{dv.margins}{0.002, 0.003 and 0.005}
\defval{dv.chosen}{0.005}
\defval{dv.m002}{0.002}
\defval{dv.pow_null}{0.949}
\defval{dv.h2}{\ensuremath{-}0.0007 [\ensuremath{-}0.0031, \ensuremath{+}0.0016]}
\defval{dv.h2_margin}{0.0028}
\defval{dv.frame_vendor}{\ensuremath{+}0.0274 [\ensuremath{+}0.0159, \ensuremath{+}0.0382]}
\defval{dv.n_sims}{1,000}
\defval{dv.sd_dev}{0.00119}
\defval{dv.sd_null}{0.00042}
\defval{dv.rc.sims}{80}
\defval{dv.rc.dev_refit}{0.00120}
\defval{dv.rc.null_refit}{0.00072}
\defval{dv.rc.null_shift}{\ensuremath{+}0.00025}
\defval{dv.rc.dev_fixed}{0.00103}
\defval{dv.rc.null_fixed}{0.00036}
\defval{dv.probe_n}{489}
\defval{dv.probe_dist}{\ensuremath{+}0.071 [\ensuremath{+}0.025, \ensuremath{+}0.122]}

\newcommand{\sbattr}{Data: StatsBomb.}
\newcommand{\Hone}{H1\xspace}\newcommand{\Htwo}{H2\xspace}\newcommand{\Hthree}{H3\xspace}\newcommand{\Hfour}{H4\xspace}
\renewcommand{\thesection}{\Alph{section}}

\def\paperID{*****}
\def\confName{CVPR}
\def\confYear{2027}

\title{What Does One Broadcast Frame Add to Expected Goals?\\A Pre-registered Comparison with Positional and Vendor xG}

\begin{document}
\maketitlesupplementary

Numbers come from three kinds of results, kept apart as in the main paper: results of the single confirmatory run (sections whose title says
``confirmatory results''), inputs frozen before the label join (pairing, kill rules, retrain gate, timing test), and
development results from models that had seen the study matches (\cref{sec:dev} only).

\section{Pre-registration and the confirmatory run}
\label{sec:prereg}

\paragraph{What was fixed.} The pre-registration (a document in the code release) fixed the hypotheses \Hone--\Hfour,
the populations and labels, the margin $\delta$, the recalibration and stacking procedure, the bootstrap, the probe
targets and thresholds, the table shells and, for \Htwo and \Hthree, the narrative for each possible reading. A hash
list pinned the analysis code (\val{pg.code_files} files), the label-free evaluation tables, the split plan, the
comparator predictions, the timing verdict and the records of two dry runs on labels permuted within matches. The dry
runs exercised the whole pipeline without real labels.

\paragraph{Decisions taken before the label join.} The principal investigator approved six decisions as recommended:
(A) the bootstrap refits recalibration and stacks in every draw; (B) one secondary row whose frame is corrupted at
decoding (it shows another moment than the shot) is excluded; (C) secondary targets without a frame stay in the table
and drop out of \Hone and \Htwo; (D) StatsBomb goals without a SoccerNet goal are excluded by a flag of the frozen
pairing; (F) the probe negative-control gate allows $R^2$ up to \val{run.gate_thr}, set after a dry run in which a
gate at zero failed on pure noise; (E) approval of the whole document. None of these decisions read a label.

\paragraph{Run.} The pre-registration record was written at 09:08:21 UTC on 2026-10-12. The confirmatory script
recomputed the hash list, compared it byte for byte with the approved one, and joined the labels once at 11:57:06 UTC;
the join record carries the pre-registration hash \val{run.prereg_sha}\ldots. The three analysis scripts exited
without error; there was no rerun and no deviation. The completion record names the pre-registration commit with a
clean tree.

\paragraph{Disclosed limits of the pre-registration.} The rules for $\delta$ and for the timing test were written
before the numbers that decided them, but in the same working session, so their precedence rests on the ticket log
rather than on separate commits. The pre-registration itself has its own commit before the label join.

\section{Population audit (confirmatory results: S1)}
\label{sec:audit}

\begin{table*}[t]
\centering\small
\begin{tabular}{@{}lrrrrrrrr@{}}
\toprule
Population & Chosen & Label-excluded & Penalties & Frame excluded & No frame & Analysed & Goals & Matches \\
\midrule
Primary & \val{cf.s1.p.chosen} & \val{cf.s1.p.label} & \val{cf.s1.p.pen} & \val{cf.s1.p.fx} & \val{cf.s1.p.nf} & \val{cf.s1.p.n} & \val{cf.s1.p.goals} & \val{cf.s1.p.matches} \\
Combined & \val{cf.s1.c.chosen} & \val{cf.s1.c.label} & \val{cf.s1.c.pen} & \val{cf.s1.c.fx} & \val{cf.s1.c.nf} & \val{cf.s1.c.n} & \val{cf.s1.c.goals} & \val{cf.s1.c.matches} \\
\bottomrule
\end{tabular}
\caption{S1, confirmatory run: populations of \Hone and \Htwo with exclusions; each row counts once, under the first
reason from the left. Combined-population frames outside the primary population are anchored at StatsBomb time plus
the half offset. \sbattr}
\label{tab:s1}
\end{table*}

\paragraph{Primary.} The paired shots lie in matches of SoccerNet's official training, validation and test splits:
\val{pa.p_train}, \val{pa.p_valid} and \val{pa.p_test}. Without the \val{cf.s1.p.pen} penalties,
\val{cf.s1.p.n} shots remain (\cref{tab:s1}).

\paragraph{Secondary.} \val{pa.secondary} StatsBomb-only shots and \val{pa.recovered} recovered goals. Their frames were
cut from \val{rd.cut} re-downloaded SoccerNet half videos, each checked against the checksum recorded when the dataset
was built; none was refused.

\paragraph{Exclusions.} \val{cf.s1.c.nf} secondary targets, all StatsBomb-only shots from \val{pg.nf_matches} matches,
lie after the end of the half video: StatsBomb counts stoppage time, and many SoccerNet videos end with the regular
half. They drop out of every rung. One StatsBomb-only shot has a frame corrupted at decoding (decision B).
\val{cf.s1.c.label} StatsBomb-only shots are StatsBomb goals without a SoccerNet goal (decision D). The
\val{cf.s1.p.pen} primary penalties and \val{pg.recovered_pen} penalties among the recovered goals are reported in
\cref{sec:pen}.

\paragraph{What the primary population misses.} Of the \val{pa.sb_unpaired} StatsBomb shots without a pair,
\val{pa.secondary} form the secondary population, \val{pa.recovered} are recovered goals, \val{pa.mispairs} lie within
\val{pa.mispair_s} of an unpaired SoccerNet shot and are dropped as possible mispairs, and two are penalties.
Conversely, \val{ti.unpaired} SoccerNet shots have no pair. \val{ti.r_shot} of them have a StatsBomb shot within
\val{pa.mispair_s}. Of the other \val{pa.no_sb_10s}, the nearest StatsBomb event of any type is another on-ball event
within \val{pa.window} for \val{ti.r_other} (usually a pass or ball receipt, so StatsBomb recorded the action as a
pass), a dead ball (throw-in, offside or foul) for \val{ti.r_dead}, and nothing within \val{pa.window} for
\val{ti.r_none}.

\section{Frozen inputs}
\label{sec:frozen}

\paragraph{Positional comparator.} The comparator learned once on \val{co.pool_shots} shots from \val{co.pool_matches}
StatsBomb matches outside the study. Its freeze-frame rung reached a pool cross-validated log-loss of \val{co.k1_ff}
against the vendor's \val{co.k1_vendor} on the same shots, a difference of \val{co.k1_gap}; K1 allowed at most
\val{co.k1_margin}, so K1 \val{co.k1_pass}. The value is the minimum over a grid of settings and therefore slightly
optimistic.

\paragraph{Retrain and gate test.} The archived frame model had \val{sp.eval_train} study matches in its training split
and \val{sp.eval_valid} in its validation split. The retrained model was compared with it on \val{gc.matches} matches
neither had seen (\cref{tab:gate}). The gate \val{gc.pass}; point estimates decide, intervals inform.

\begin{table*}[t]
\centering\small
\begin{tabular}{@{}lcccc@{}}
\toprule
Metric & Retrained & Archived & Difference [\val{pg.ci95}] & Tolerance \\
\midrule
Log-loss & \val{gc.ll_new} & \val{gc.ll_old} & \val{gc.ll_diff} & \val{gc.tol_ll} \\
AUC & \val{gc.auc_new} & \val{gc.auc_old} & \val{gc.auc_diff} & \val{gc.tol_auc} \\
\bottomrule
\end{tabular}
\caption{Retrain gate on \val{gc.shots} shots, frozen before the label join.}
\label{tab:gate}
\end{table*}

\paragraph{Timing-source test.} \label{sec:timing}The rule was written before the comparison. Substitutions, which
SoccerNet marks when the broadcast shows them, decide between a shared clock and two independent annotations
(\cref{tab:timing}). The verdict is \val{ti.verdict}. A copied clock would not give substitutions offset by seconds.
The corner margin over its threshold was thin, \val{pg.corner_margin}. Two consequences follow: the timing-residual
arm of \Hfour runs, always with the restriction to $\pm200$\,ms as a sensitivity; and secondary frames, cut at
StatsBomb time, can lie a few hundred milliseconds from the contact.

\begin{table}[t]
\centering\small
\setlength{\tabcolsep}{3pt}
\begin{tabularx}{\linewidth}{@{}>{\raggedright\arraybackslash}Xrl@{}}
\toprule
Quantity & Value & Rule \\
\midrule
Substitutions in the rule set & \val{ti.n_subs} & at least \val{ti.min_subs} \\
Substitutions within \val{ti.agree_ms} & \val{ti.subs_share} & from \val{ti.shared}: shared clock \\
Median $|r|$ of substitutions & \val{ti.subs_median} & from \val{ti.lag}: independent \\
Shots within $\pm200$\,ms, leave-one-out offset & \val{ti.shots} & at least \val{ti.kick} \\
Corners within $\pm200$\,ms & \val{ti.corners} & at least \val{ti.kick} \\
\bottomrule
\end{tabularx}
\caption{Timing-source test, frozen before the label join; $r$ is the SoccerNet time minus the StatsBomb time after the
half offset. \sbattr}
\label{tab:timing}
\end{table}

\section{Additional confirmatory results}
\label{sec:moreresults}

\subsection{\Hone: key contrasts (descriptive)}

\begin{table*}[t]
\centering\small
\begin{tabular}{@{}lcc@{}}
\toprule
Contrast (A $-$ B) & Primary & Combined \\
\midrule
Frame $-$ vendor & \val{cf.c.p.frame_vendor} & \val{cf.c.c.frame_vendor} \\
Frame $-$ positions & \val{cf.c.p.frame_ff} & \val{cf.c.c.frame_ff} \\
Frame $-$ dist.\ + angle & \val{cf.c.p.frame_da} & \val{cf.c.c.frame_da} \\
Vendor $-$ positions & \val{cf.c.p.vendor_ff} & \val{cf.c.c.vendor_ff} \\
Positions $-$ events & \val{cf.c.p.ff_event} & \val{cf.c.c.ff_event} \\
Events $-$ dist.\ + angle & \val{cf.c.p.event_da} & \val{cf.c.c.event_da} \\
Dist.\ + angle $-$ constant & \val{cf.c.p.da_const} & \val{cf.c.c.da_const} \\
\bottomrule
\end{tabular}
\caption{\Hone, confirmatory run, descriptive: log-loss differences of recalibrated rungs with \val{pg.ci95}
intervals. Combined-population frames outside the primary population are anchored at StatsBomb time plus the half
offset. \sbattr}
\label{tab:contrasts}
\end{table*}

\Cref{tab:contrasts} gives the key contrasts of both populations. The largest step is from the constant to distance
and angle, the smallest from events to positions. Full tables, with Brier score, calibration slope and intercept and
the raw variants, are in the release.

\subsection{\Hthree: descriptive targets}
\label{sec:h3more}

\begin{table*}[t]
\centering\footnotesize
\setlength{\tabcolsep}{3pt}
\begin{tabular}{@{}llllll@{}}
\toprule
Target & Metric & G-combo fold 0 [\val{pg.ci95}] & Reading, fold 0 & Reading, fold 1 & G $-$ DINOv2, fold 0 / fold 1 \\
\midrule
Angle & $R^2$ & \val{cf.h3.f0.angle} & encodes & encodes & G $>$ DINOv2 / G $>$ DINOv2 \\
Defenders in cone (count) & $R^2$ & \val{cf.h3.f0.ff_def_cone} & encodes & encodes $\dagger$ & G $>$ DINOv2 / no difference \\
At least two defenders in cone & AUC & \val{cf.h3.f0.ge2} & encodes & encodes & G $>$ DINOv2 / no difference \\
Goalkeeper depth & $R^2$ & \val{cf.h3.f0.gk_depth} & encodes $\dagger$ & undecided $\dagger$ & G $>$ DINOv2 / G $>$ DINOv2 \\
Goalkeeper lateral offset & $R^2$ & \val{cf.h3.f0.gk_off} & does not encode & does not encode & G $>$ DINOv2 / G $>$ DINOv2 \\
Goalkeeper distance to goal & $R^2$ & \val{cf.h3.f0.gk_dist} & undecided $\dagger$ & undecided & G $>$ DINOv2 / no difference \\
Opponents near the shooter$^\ast$ & $R^2$ & \val{cf.h3.f0.opp} & does not encode & does not encode & G $>$ DINOv2 / G $>$ DINOv2 \\
Header & AUC & \val{cf.h3.f0.head} & encodes & encodes & no difference / G $>$ DINOv2 \\
\bottomrule
\end{tabular}
\caption{\Hthree, confirmatory run, descriptive targets. Thresholds: $R^2$ \val{cf.h3.thr.angle}, AUC
\val{cf.h3.thr.head}; readings follow the lower bound as for the confirmatory targets. $\dagger$: lower bound within
\val{pg.near} of the threshold. ``No difference'': the interval of G $-$ DINOv2 includes zero. $^\ast$Opponents within three StatsBomb
pitch units (yards) of the shooter. \sbattr}
\label{tab:h3desc}
\end{table*}

\Cref{tab:h3desc} lists the eight descriptive targets. The fold models differ only on descriptive targets and on
difference readings; the readings of the confirmatory targets agree. The descriptive thresholds were set with the
prototype's development numbers known. Probe intervals are conditional: the draws reweight fixed out-of-fold
predictions and the probes are not refitted, so they omit the variability of the regularisation choice and the
coefficients and are narrower than the full uncertainty.

\subsection{\Hfour: leakage diagnostic (descriptive)}
\label{sec:h4more}

\begin{table*}[t]
\centering\small
\begin{tabular}{@{}lcc@{}}
\toprule
Model, with controls & All rows & Within $\pm200$\,ms \\
\midrule
Slope of the gap on $r$ [logit/s] & \val{cf.h4.slope} & \val{cf.h4.slope_200} \\
Anchor more than 40\,ms before StatsBomb & \val{cf.h4.early} & \val{cf.h4.early_200} \\
Anchor more than 40\,ms after StatsBomb & \val{cf.h4.late} & \val{cf.h4.late_200} \\
Shots & \val{cf.h4.n} & \val{cf.h4.n200} \\
\bottomrule
\end{tabular}
\caption{\Hfour, confirmatory run, descriptive: timing-residual arm. The gap is logit(frame) $-$ logit(vendor); $r$ is
the timing residual. \sbattr}
\label{tab:h4}
\end{table*}

The offset contrast G(0) $-$ G($-400$\,ms) on on-target shots is \val{cf.h4.off_on} with controls and
\val{cf.h4.off_on_u} without. A positive contrast on on-target shots has the two explanations of the on-target
coefficient and a third: play develops during the last $400$\,ms before the anchor. Over all rows the interval of the
residual slope lies above zero; within $\pm200$\,ms it is wide and includes zero (\cref{tab:h4}).

\subsection{Penalties}
\label{sec:pen}

Penalties are outside \Hone--\Hfour. The frame model does not score them: it gives every penalty the pool penalty rate
of its calibration. \Cref{tab:pen} gives raw probabilities, without intervals.

\begin{table*}[t]
\centering\small
\begin{tabular}{@{}lrrcccccc@{}}
\toprule
 & & & \multicolumn{2}{c}{Vendor} & \multicolumn{2}{c}{Frame} & \multicolumn{2}{c}{Positions} \\
\cmidrule(lr){4-5}\cmidrule(lr){6-7}\cmidrule(lr){8-9}
 & $n$ & Goals & mean $p$ & LL & mean $p$ & LL & mean $p$ & LL \\
\midrule
Primary & \val{cf.pen.p.n} & \val{cf.pen.p.goals} & \val{cf.pen.p.vp} & \val{cf.pen.p.vll} & \val{cf.pen.p.fp} & \val{cf.pen.p.fll} & \val{cf.pen.p.pp} & \val{cf.pen.p.pll} \\
Combined & \val{cf.pen.c.n} & \val{cf.pen.c.goals} & \val{cf.pen.c.vp} & \val{cf.pen.c.vll} & \val{cf.pen.c.fp} & \val{cf.pen.c.fll} & \val{cf.pen.c.pp} & \val{cf.pen.c.pll} \\
\bottomrule
\end{tabular}
\caption{Penalties, confirmatory run, descriptive (S9). \sbattr}
\label{tab:pen}
\end{table*}

\section{Context from models that had seen these matches (development results)}
\label{sec:dev}

These numbers are not results of the study. The archived frame model or the prototype computed them, and the archived
model had trained on most of these matches. They explain choices of the pre-registration.

\paragraph{Choice of $\delta$.} The rule, written before the power simulation, took the smallest candidate margin with
power of at least \val{dv.target} under the development truth. TOST power was \val{dv.pow002}, \val{dv.pow003} and
\val{dv.pow005} at margins \val{dv.margins}, so the rule chose $\delta = \val{dv.chosen}$. The development truth
assumed a frame gain about twice the one that \Htwo showed on the development table. Had the frame added nothing,
the margin \val{dv.m002} would already have had power \val{dv.pow_null}. The rule was kept unchanged. Hence
``equivalent at $\delta$'' is a weak claim, and the smallest equivalence margin is reported next to it.

\paragraph{\Htwo on the development table.} The archived model gave \val{dv.h2} with smallest margin
\val{dv.h2_margin}; frame $-$ vendor was \val{dv.frame_vendor}.

\paragraph{Bootstrap with refitting.} In the power simulation (\val{dv.n_sims} studies per scenario) the sampling SD
of the \Htwo difference was \val{dv.sd_dev} under the development truth and \val{dv.sd_null} under the null truth. A
check on \val{dv.rc.sims} simulated studies gave a mean SD of the refitting bootstrap of \val{dv.rc.dev_refit} under the
development truth, in line with the spread, and \val{dv.rc.null_refit} under the null truth, clearly too wide; there
the draws lay on average \val{dv.rc.null_shift} above the estimate. A bootstrap with fixed predictions gave
\val{dv.rc.dev_fixed} and \val{dv.rc.null_fixed}, too narrow in both scenarios. The refitting bootstrap was chosen
because a too-wide interval only makes equivalence harder to show, whereas a too-narrow one would make TOST liberal.

\paragraph{Probes on single fold models.} The prototype mixed the embeddings of five fold models in one probe and saw
no advantage of the fine-tuned model over DINOv2. A probe on the single fold-0 model, on \val{dv.probe_n} shots, gave
a distance difference of \val{dv.probe_dist}. The pre-registration therefore wrote both fine-tuning narratives
neutrally.

\section{Release package}
\label{sec:release}

The package ships five tables and the code archive: the match map (every SoccerNet game with its StatsBomb match
identifier or the reason it has none), the pair table (every index shot of the paired games with its StatsBomb event
identifier), the recovered goals, the frame model's calibrated logits and probabilities at offsets 0 and $-400$\,ms for
every evaluation row, and the split plan of the retrain. A manifest lists each column with its provenance (SoccerNet
identifier, StatsBomb identifier, or ours) and each source column left out with the reason. Left out are: every
StatsBomb value, the comparator's predictions (functions of StatsBomb data only; the manifest gives the checksum a
rebuild must reach), half offsets and timing residuals (with the SoccerNet time they give back the StatsBomb
timestamp), SoccerNet annotation content, embeddings, probe targets, frames, clips and weights. Frame predictions are
not penalty-routed, because the penalty flag is a SoccerNet annotation or a StatsBomb value.

\section{Number check}
\label{sec:numbercheck}

No number in the paper or this supplement is typed by hand. Each is a macro filled from a claim that names a number of
the study report, which reads it from a result file by a fixed locator (CSV row and column, JSON path, or a quotation
of the pre-registration). A script renders every claim, rewrites nothing silently and fails when the paper's macros,
the report's numbers, the figures or the SHA-256 of any result file disagree with the sources, when a number in the
text is not a macro, or when a confirmatory or development number stands outside a section labelled as such.

\end{document}
